

The Kitaev model \cite{Kitaev2006} is an exactly solvable spin-half model that exhibits exchange frustration, topological order, long-range entanglement, and fractionalization into non-Abelian anyons and  Majorana fermions.
Beyond the theoretical fascination of a simple and very rich model, the interest in Kitaev physics goes to its possible experimental realization.


This chapter introduces the Kitaev model, the exact solution using Majorana fermions, and the topological property from bulk and edge perspectives.

The low energy spin excitations about the ground state fractionalize into Majorana fermions and $\mathbb{Z}_2$ vortices with a Majorana zero mode localized on it. Depending on the anisotropy of the spin interactions, the ground state can be in a gapped or a gapless phase.
The gapped phase is connected to the toric code.
The gapless phase becomes gapped under a magnetic field, which gives rise to topological non-trivial properties such as non-zero Chern number and non-Abelian anyonic statistics.




\section{The spin Hamiltonian}
In this model, we have spin-half degrees of freedom $\mb{S}_i=\frac{1}{2}\bm{\sigma}_i$ living in a honeycomb lattice.
The Pauli operators satisfy the algebra
\begin{equation}
    \sigma_i^\alpha \sigma_i^\beta = \mathbb{1} \, \delta^{\alpha\beta} + \im \sum_\gamma \varepsilon^{\alpha\beta\gamma} \sigma_i^\gamma \; , \quad [\sigma_i^\alpha ,  \sigma_j^\beta] = 0 \; ,
\end{equation}
for $\alpha,\beta,\gamma \in \{x,y,z\}$.
The position vector $\mb{R}_{i}$ for the site $i=(r_{i1},r_{i2},\nu_i)$ is indexed by the unit cell  $(r_{1},r_{2}) \in \mathbb{Z}_{L_1}\times\mathbb{Z}_{L_2}$ and the sublattice index $\nu=0,1=\rmA,\rmB$. We use periodic boundary conditions in the geometry of $L_1\times L_2$ unit cells depicted in Fig.~\ref{fig:basis_vec}.

The Kitaev Hamiltonian with the $\kappa$ term is
\begin{equation}
    H = \sum_{\alpha}
    \sum_{\braket{ij}_\alpha} K_\alpha \sigma_i^\alpha\sigma_j^\alpha
    - \kappa \sum_{\gamma (\beta \alpha)}
    \sum_{\braket{ij}_\alpha\braket{kj}_\beta}   \sigma_i^\alpha \sigma_j^\gamma  \sigma_k^\beta
    \; , \label{eq:H0spin}
\end{equation}
where $\gamma$ is determined so that $(\gamma\beta\alpha)$ is a cyclic permutation of $(xyz)$ and $\braket{ij}_\alpha$ is a sum over all bonds $(ij)$ that are in a $\alpha$ bonds. In the sum over bonds, we count each bond once, so if $(ij)$ is in the sum, then we do not count $(ji)$. This allows us to fix the first site to be in the A sublattice and the second in the B sublattice.
Usually, the Kitaev model refers only to the first term  ($\kappa=0$). %, however in Ref.\cite{Kitaev2006} bor.
The second term is the most relevant perturbation generated by the applyed magnetic field in the Kitaev Hamiltonian.



%Let $\bm{\delta}_{\alpha}$ be the vectors connecting nearest neighbors (NN) in $\alpha$-bonds and let $\mathbf{N}_{\alpha}$ be the vectors connecting next nearest neighbors (NNN) ``complementary'' to $\alpha$-bond, these vectors are defined in Fig.~\ref{fig:basis_vec}. Note that the NNN bond $\gamma$, cyclically complementary to $(\alpha\beta)$, is obtained by\footnote{$\gamma(\beta\alpha)$ is cyclic equivalent to $(\alpha)\gamma(\beta)$,  e.g. $\gamma=y \to (\alpha,\beta)=(x,z)$. } $\mathbf{N}_{\gamma(\beta\alpha)}=\bm{\delta}_{\alpha} -  \bm{\delta}_{\beta}$.


\begin{figure}[h!]
    \centering
    \includegraphics[width=.95\textwidth]{Figures/2/lattice-basis-vec.pdf}
    \caption{ Honeycomb lattice with basis vectors $\mathbf{a}_{1/2}$,  nearest neighbors $\bm{\delta}_\gamma$, and next nearest neighbors $\mathbf{N}_\gamma$.
    The unit cell is chosen along $z$-bonds with sublattices $\rmA$ (black circles) and $\rmB$ (white circles).  The color code used here is \textcolor{darkred}{Red}, \textcolor{darkgreen}{Green}, and \textcolor{darkblue}{Blue} for \textcolor{darkred}{$x$-}, \textcolor{darkgreen}{$y$-}, and \textcolor{darkblue}{$z$-}bonds and also for next nearest neighbor `bonds' in dotted lines.
    }
    \label{fig:basis_vec}
\end{figure}

%\clearpage


We can rewrite the sum in the spin Hamiltonian eq.~\eqref{eq:H0spin} over bonds in the honeycomb lattice as a sum over one sublattice
\begin{equation}
    H =
    \sum_{i\in \mathrm{A}}\sum_{\alpha}  K_\alpha \sigma_i^\alpha\sigma_{i_\alpha}^\alpha
    - \kappa
    \sum_{i\in \mathrm{A}} \sum_{\gamma(\beta \alpha )}   \sigma_i^\alpha \sigma_{i_{\alpha}}^\gamma  \sigma_{i_{\alpha\bar{\beta}}}^\beta
     - \kappa
    \sum_{i\in \mathrm{B}} \sum_{\gamma(\beta \alpha )}   \sigma_i^\alpha \sigma_{i_{\bar{\alpha}}}^\gamma  \sigma_{i_{\bar{\alpha}{\beta}}}^\beta
    \; ,  \label{eq:2-kitaev}
\end{equation}
where the position vector for the nearest neighbors and next-nearest neighbors to the site $i$ at $\mathbf{R}_i$ are
\begin{align}
    \mb{R}_{i_\alpha\;} &= \mathbf{R}_i + \bm{\delta}_{\alpha} \; ,  \\[5pt]
    \mb{R}_{i_{\alpha\bar{\beta}}} &= \mathbf{R}_i + \bm{\delta}_{\alpha} -  \bm{\delta}_{\beta}  = \mathbf{R}_i + \mathbf{N}_{\gamma} \; , \quad (\alpha\gamma\beta)\overset{\scriptsize\text{cyclic}}\sim(xyz)
\end{align}
 and the bar represents traversing the bond in the reversal direction $\mb{R}_{i_{\bar{\alpha} }} = \mathbf{R}_i  -\bm{\delta}_{\alpha}$.
% with  $\mathbf{R}_i=R_{i,1} \mb{a}_1 + R_{i,2} \mb{a}_2$ and $R_{i,1},R_{i,2}\in\mathbb{Z}$.
Note that the sum in $\gamma$ only has $3$ terms, i.e for $\gamma=x,y,z$ we have respectively $(\alpha,\beta)=(z,x),(x,z), (y,x)$.  The basis vectors $\mathbf{a}_{1/2}$, and the (next-) nearest vectors $\bm{\delta}_\gamma$ and $\mb{N}_\gamma$  are defined in Fig.~\ref{fig:basis_vec}.




The Kitaev Hamiltonian has an extensive number of integrals of motion $W_p$ defined for each hexagonal plaquette in the honeycomb lattice. The hexagons in the honeycomb lattice are in 1-to-1 correspondence to the unit cells, so we can equivalently index each plaquette $p$ with a site $i$ in the $\rmA$ sublattice. The hexagonal plaquette $p$ is the regular polygonal region with vertices $(p_1, p_2, ...,p_6)=\left( i, i_{\bar x},i_{\bar x z}, i_{\bar x z \bar y} , i_{\bar y z}, i_{\bar y} \right)$ in which the edges $\partial p = (p_1p_2,p_2p_3,...,p_6p_1)$ are respectively in bonds $(x,z,y,x,z,y)$ as depicted in figure \ref{fig:2-plaquette}.
For each plaquette $p$, we define the observable
\begin{equation}
\begin{split}
    W_p
    &=
    \sigma_{p_1}^z \sigma_{p_2}^y \sigma_{p_3}^x \sigma_{p_4}^z \sigma_{p_5}^y \sigma_{p_6}^x
    \\[6pt]
    &=
    (\sigma_{p_1}^x \sigma_{p_2}^x)
    (\sigma_{p_2}^z \sigma_{p_3}^z)
    (\sigma_{p_3}^y \sigma_{p_4}^y)
    (\sigma_{p_4}^x \sigma_{p_5}^x)
    (\sigma_{p_5}^z \sigma_{p_6}^z )
    (\sigma_{p_6}^y \sigma_{p_1}^y )
\end{split}
\end{equation}
which can either be written as a product over sites $\prod_{i\in p}\sigma_i^{\alpha_i}$ with the spin component $\alpha_i$ pointing outwards of the plaquette or be written as a product over bonds, see Fig.~\ref{fig:2-plaquette}.
The former form makes it clear that the operators $\sigma_i^{\alpha_i}$, for components $\alpha$ exterior to the plaquette $p$, commutes with $W_p$ while $\sigma_i^{\gamma}$ anticommutes for $\gamma\neq \alpha_i$. This anticommutativity implies that the operator $\sigma_i^{\gamma}$ flips the sign of $W_p$ for two hexagons that share the bond\footnote{For the bonds $\braket{ij}_\gamma$ in the honeycomb lattice, it is enough to give one site, e.g. $i$, and the type of bond, $\gamma$, to uniquely determine the bond $(ij)$. } $\braket{ij}_\gamma$,
\begin{equation}
    \sigma_i^{\gamma} \, W_p = -  W_p \, \sigma_i^{\gamma} \; , \quad \braket{ij}_\gamma \in \partial p \; . \label{eq:flux}
\end{equation}
While, from the form $W_p = \prod\limits_{\braket{ij}_{\alpha} \in \partial p} (\sigma_{i}^\alpha \sigma_{j}^\alpha)$, we can see that the plaquette operators commute with the $K$-term in the Kitaev Hamiltonian since it is a sum of mutually commuting terms form $\sigma_{i}^\alpha \sigma_{j}^\alpha$ for different bonds $\braket{ij}_\alpha$.  In the following I will call these terms ``bond operators''.  The $\kappa$ term also commutes with $W_p$, which follows by writing it as
\begin{equation}
    \sigma_i^\alpha \sigma_{i_{\bar{\alpha}}}^\gamma  \sigma_{i_{\bar{\alpha}{\beta}}}^\beta
    =  \im
    (\sigma_i^\alpha
    \sigma_{i_{\bar{\alpha}}}^\alpha)
    (\sigma_{i_{\bar{\alpha}}}^\beta
    \sigma_{i_{\bar{\alpha}{\beta}}}^\beta)
\end{equation}
which are also products of the bond operators for bonds $\alpha$ and $\beta$ (with $\alpha\gamma\beta$ in a cyclic permutation of $xyz$) .
The factorization in a product of bond operators is what guarantees the exact solvability of the model.



\begin{figure}[h!]
    \centering\vspace*{-3mm}%
    \raisebox{0.22\height}{\includegraphics[width=0.3\textwidth]{Figures/2/plaquette-p.pdf}}\hfill%
    \includegraphics[width=0.58\textwidth]{Figures/2/2vortices.pdf}
    \caption{Vertices for a hexagonal plaquette $p$. Configuration with two adjacent vortices $W_p=-1$ created by acting with $\sigma_i^z$ or $\sigma_j^z$ on the flux-free configuration. }
    \label{fig:2-plaquette}
\end{figure}

The Kitaev Hamiltonian commutes with the set of mutually commuting observables $W_p$.
To find the eigenstates of the Hamiltonian $H$ we can reduce to search within each eigenspace of $\{W_p\}$, which are invariant subspaces of the Hamiltonian.
Each of these subspaces is said to be a flux sector. The flux of a plaquette is given by the eigenvalue $w_p$ of $W_p$ in a given state of the flux sector. Since all operators $W_p$ are Hermitian and satisfy $(W_p)^2=1$ we must have two orthogonal eigenvectors per plaquette with corresponding eigenvalues $w_p=\pm 1$. As these fluxes take binary value, we said that $W_p$ measures the $\mathbb{Z}_2$ flux. For the sector where $w_p=e^{\im \pi} = -1$, we said that we have a $\pi$-flux or a $\mathbb{Z}_2$ vortex in the plaquette $p$. % On the nomenclature: In the Kitaev model, we will see that for a (sufficiently) isolated vortex the system traps a Majorana zero mode (MZM),  the resulting bound state of a vortex and a MZM is called vison and has non-Abelian anyonic properties.

As discussed before eq.~\eqref{eq:flux}, the spin operator $\sigma_i^\gamma$ flips (changes the sign) of the flux in the two plaquettes $W_p$. Then, on top of a flux-free region, the spin operator creates a pair of vortices. A successive application of the spin operator on the same bond annihilates the two vortices. However, by applying the spin operator on another bond of these hexagons, it will move the corresponding vortex. As a consequence of this, any physical operation of applying a finite number of spin operators can not change the parity for the number of vortices. Out of the flux-free configuration, one can only arrive at configurations with an even number of vortices.

%We replace the flux operator $W_p$ operators by their eigenvalues $W_p \to w_p$.

%The Kitaev Hamiltonian is block diagonal within each eigenspace.
In the basis of $W_p$, the Kitaev Hamiltonian is block diagonal.
%The replacement of $W_p \to w_p$ in the Hamiltonian is not straightforward to see in the spin language.
This simplifies the problem as we can restrict to each flux sector at a time.
However, there is the immediate problem of determining which flux configuration to choose. For $N$ unit cells, we expect to have exponentially many ($\sim2^N$)  possible flux configurations. Fortunately, due to a theorem by Lieb \cite{Lieb_1994}, it is known that the ground state is in a flux-free configuration for all hexagonal plaquettes.
Here we will be interested in the ground state and in the configuration with few isolated vortices.


%The flux operator $W_p$, measures the $\mathbb{Z}_2$-flux in the plaquette.

In characterizing the flux configurations, we can use the Wilson loop operator $W(\ell)$ that measures the flux through the loop $\ell$,
\begin{equation}
    W(\ell) = \prod_{\braket{ij}_\alpha \in \ell } \sigma_i^\alpha \sigma_j^\alpha \; ,
\end{equation}
where the sum is over the edges $(\ell_n,\ell_{n+1}
)$ (bonds) in a closed loop $\ell=\{\ell_1,\ell_2,...,\ell_{L},\ell_1\}$.
All these operators square to $1$ and commutes with the Hamiltonian.
For $\ell$ equal to the vertices of the plaquette $p$, the Wilson loop is the flux operator $W(p)=W_p$. Moreover, it satisfies a discrete Stoke's theorem property that the flux on a loop is the product of the flux over plaquettes inside the loop:
\begin{equation}
    W(\ell) = \prod_{p \in \mathcal{S}} W_p \; ; \quad \partial \mathcal{S} = \ell  \label{stokes}
\end{equation}
Note that the orientation of the loop $\ell$ (hence which surface $\mathcal{S}$ with boundary $\ell$) does not matter because of the global constraint
\begin{equation}
    \prod_p W_p = +1 \; ,
\end{equation}
which implies that the flux inside the loop equals the flux outside $W(\ell)=W(\bar\ell)$. The product structure of  \eqref{stokes} implies that its value $ W(\ell)$ is not altered if we deform the loop $\ell$ in a flux-free region. This allows us to contract loops as long as we do not cross the position of a vortex. In analogy to contour integration, $W(\ell)$ is the line integral over $\ell$ and $\mathbb{Z}_2$ vortices are singularities.

On a plane, all loops $\ell$ are contractible, and thus all Wilson loop operators are trivial in the flux-free configuration $W_p=+1$. This implies that we only have one ground state.
The situation changes for the realization of the Kitaev model in surfaces with higher genus $g$, where one finds degenerate\footnote{In a finite system they have nearly degenerated energies  $\sim \exp(-|\ell|/\xi) $ with $|\ell|$ the length of the loop in the surface and $\xi$ the coherence length that goes with the inverse of the gap to the excited states.} ground states.
Here we are interested in the realization on the torus, which is obtained by imposing periodic boundary conditions on \eqref{eq:H0spin}.


\begin{figure}[t]
    \centering
    \includegraphics[width=.3\textwidth]{Figures/2/W-loop-2.pdf}\;\;%
    \includegraphics[width=.3\textwidth]{Figures/2/T-op3.pdf}
    \caption{(a) Wilson operators on the non-trivial loops $\ell_{1,2}$ along directions $L_{1/2}$ of the torus. (b) Operators that wrap the vortices along these loops.
    The operators $W$ and $\mathcal{T}$ commute the ones defined on the same loop, and anticommute for different $W^{(1)}\mathcal{T}_2 W^{(1)}=-\mathcal{T}_2$ and $W^{(2)}\mathcal{T}_1 W^{(2)}=-\mathcal{T}_1$.}
    \label{fig:W-T}
\end{figure}



The torus has two non-contractible loops, say $\ell_1$ and $\ell_2$, which circulate in their poloidal and toroidal directions as in Fig.~\ref{fig:W-T}(a). For any such loops, we define the fluxes
\begin{equation}
    W^{(1)}=W(\ell_1) \, , \quad
    W^{(2)}=W(\ell_2) \; ,
\end{equation}
which measures the $\mathbb{Z}_2$-flux through the %\footnote{E.g. introducing magnetic field in a cylinder}
``holes'' of the torus.  Traversing each of the loop $\ell_{1/2}$ twice will give the identity, but   $W^{(1)}W^{(2)}$ is non-trivial. Thus, in the torus, there are  $4$ states in the vortex-free configurations indexed by\footnote{Also indexed by $W_p=+ \; \forall p$ for the vortex-free configurations. In each of these sectors, there are many states, and within each sector, these are the ones with the lowest energy. } the eigenvalues of $(W^{(1)},W^{(2)})=(+,+),(-,+),(+,-),(-,-)$. %Lets denote them by $\ket{\mathbb{1}},\ket{\sigma_x},\ket{\sigma_y},\ket{\sigma_x\sigma_y}$, for which $\sigma$ is associated with the $\mathbb{Z}_2$-flux through the corresponding hole of the torus, and $\mathbb{1}$ the vaccum.
These states may be obtained from $\ket{(+,+)}$ by the action of $\mathcal{T}_{1/2}$. Which creates a pair of $\mathbb{Z}_2$ fluxes in adjacent plaquettes and wraps one of these particles across the corresponding loop $\ell_{1/2}$ ending up annihilating the two in the original position.  A specific construction of these operators is depicted in Fig.~\ref{fig:W-T} in which each line corresponds to the spin operator $\sigma_i^\alpha$ in one of the sites of the bond (e.g. with $i$ in the $\rmA$ sublattice). The action of $\mathcal{T}_{1}$ introduce a flux along the loop $\ell_2$, and $\mathcal{T}_{2}$ on $\ell_1$, so that\vspace*{-2mm}%
\begin{equation}%
\ket{(-,+) }\propto \mathcal{T}_2\ket{(+,+)}, \;
\ket{(+,-) }\propto \mathcal{T}_1\ket{(+,+)}, \;
\ket{(-,-) }\propto \mathcal{T}_1\mathcal{T}_2\ket{(+,+)} \; .\vspace*{-1mm}%
\end{equation}
This construction is very general and only depends on the genus of the surface. However, even though all these four states have zero flux through each hexagonal plaquette, this does not mean that they are eigenstates of the Kitaev Hamiltonian with the same energy.
%The action of $\mathcal{T}$ on the Hamiltonian is not a symmetry, but to flip the sign of the exchanges along the bonds shown in Fig.~\ref{fig:W-T}(b).
It turns out that in the Abelian phase of the Kitaev model, these four states are degenerate  (the same as in the Toric code \cite{Kitaev2003}).
However, in the non-Abelian phase, there are only three ground states, due to the fact that non-commutativity of $\mathcal{T}_{1/2}$ creates an excitation.

To further proceed in understanding the Kitaev model, we will use a spin representation that allows us to analytically solve the Kitaev model.









\section{Fermionic model}

The solution proposed by Kitaev in \cite{Kitaev2006} invokes the representation of spins in terms of fermionic operators.
The trick of using the so-called partons (fermions or bosons) operators to represent the spin has been known for a long time, for instance, Ref. \cite{Abrikosov1965, Baskaran_1987}.
However, the major downside of these representations is that one needs to consider an enlarged Hilbert space with a set of holonomic constraints.
A particularity of Kitaev's model is that we can deal with the enlarged space and find exact solutions using this representation.
To the readers unfamiliar with Majorana fermions and parton representations the appendix \ref{App-Parton} is recommended.

The spin representation used by Kitaev trades each spin-half $\bm{\sigma}_i$ by four Majorana fermions $(c^0_i,c^1_i,c^2_i,c^3_i)$.
This enlarges the local Hilbert space from\footnote{The Hilbert space for two Majorana fermions has dimension $2$ (or dimension $\sqrt{2}$ for each Majorana).} $2$ to $(\sqrt{2})^4=4$.
The spin representation is
\begin{equation}
 \sigma_i^\alpha = \im c^0_i c^\alpha_i\, , \quad \rmD_i = c^0_ic^1_ic^2_ic^3_i=+1 \quad \forall i \; ,
\end{equation}
with Majorana fermions $c^\mu$ being Hermitian operators normalized so that $\{c^\mu_i,c^\nu_j\} = 2 \delta^{\mu\nu} \delta_{ij}$.
The spin operators commute for different sites and by noting that $\rmD_i^2=\mathbb{1}$ and
$  c^0_i c^\gamma_i = -\frac{1}{2} \epsilon^{\alpha\beta\gamma} \rmD_i  c^\alpha_i c^\beta_i $
one can easily check that the Pauli operators
 satisfy the local algebra
\begin{equation}
    \sigma_i^\alpha \sigma_i^\beta  = \delta^{\alpha\beta} \mathbb{1}   +  \rmD_i \; \im \epsilon^{\alpha\beta\gamma} \sigma_i^\gamma \; ,
\end{equation}
which is the spin-half algebra in the subspace where $\rmD_i=+1$. Note that any bilinear $c_i^\mu c_i^\nu$, such as $\sigma_i^\alpha$, commutes with the operator $\rmD_i  = c^0_ic^1_ic^2_ic^3_i $. This means that we never leave the subspace where $\rmD_i=1 \; \forall i$ (called the physical space) by applying physical operators such as the spin operator $\frac{1}{2}\bm{\sigma}_i$ or the Hamiltonian.

The operator $\rmD_i$ implements a $\mathbb{Z}_2$ gauge transformation
\begin{equation}
    c_j^\mu \quad \longrightarrow \quad \rmD_i \,c_j^\mu\, \rmD_i = %\sigma_{ij}\; c_j^\mu \; \quad \sigma_{ij}=
    \left\{\begin{array}{cc}   -c_j^\mu & i=j \\[6pt]   +c_j^\mu & i\neq j \end{array}\right. \;.
\end{equation}
It is a redundancy of the representation as this symmetry transformation keeps all observables, such as the spin and the flux, invariant.


The constraint $\rmD_i = c^0_ic^1_ic^2_ic^3_i=+1$  can be presented in bilinear form in the fermions as
\begin{equation}\left.
    \begin{array}{cc}
        c^0_i c^1_i &= - c^2_i c^3_i\\
        c^0_i c^y_i &= - c^3_i c^1_i \\
        c^0_i c^z_i &= - c^1_i c^2_i
    \end{array}\right\} \quad \Rightarrow \quad
        c^0_i c^\gamma_i = - c^\alpha_i c^\beta_i   \; .
\end{equation}
Here, for fixed $\gamma$, I am using the short-hand notation $c^\alpha_i c^\beta_i = -\frac{1}{2} \epsilon^{\alpha\beta\gamma} c^\alpha_i c^\beta_i$ with implicit sum over $\alpha$ and $\beta$.

Using this representation one gets a quartic term in $c^\mu$ in the Kitaev Hamiltonian ($\kappa=0$). However, the exchange interactions in the Kitaev Hamiltonian are so that the $c^\alpha$ terms, for $\alpha=1,2,3$, only appear in the combination $\im c_i^\alpha c_j^\alpha$ in an $\alpha$-bond $\braket{ij}_\alpha$. This greatly simplifies the problems, as the operators
\begin{equation}
    \hat{u}_{\braket{ij}_{\alpha} }= \im c_i^\alpha c_j^\alpha
\end{equation} commute with themselves for different bonds  $\braket{ij}_\alpha$, giving a constant of motion.
The notation $\braket{ij}_{\alpha}$ refers to a specific ordered $\gamma$-bond with $i$ in the $\rmA$ sublattice, which automatically fix $j$. The gauge fields are defined for each bond in the honeycomb lattice.
Since these operators are Hermitian and $(\hat u_{\braket{ij}})^2=\mathbb{1}$,  there has two orthogonal eigenvectors per bond with eigenvalues $u_{\braket{ij}} = \pm 1$.
These constants of motion in the Kitaev model form a $\mathbb{Z}_2$-gauge field $u_{\braket{ij}}$.
So that we can write the Hamiltonian as a quadratic term in $c^0$ times the gauge field.
For the $\kappa$-term in eq. \eqref{eq:2-kitaev} we expect a generic sixth term in $c^\mu$ but it can also be written as a quadratic term in $c^0$, but with a prefactor quadratic in $\hat{u}_{\braket{ij}}$.


The common eigenstate of $\hat{u}_{\braket{ij}_{\alpha} }$ form a basis for a vector space  $\mathcal{H}_u$ subspace of the extended Hilbert space $\mathcal{H}_f$.
The states are $\ket{u} \in \mathcal{H}_u$ and satisfies the eigenvalue equation
\begin{equation}
    \hat{u}_{\braket{ij}_{\alpha} }\ket{u}={u}_{\braket{ij}_{\alpha} }\ket{u} \; ,
\end{equation}
which reads the gauge field ${u}_{\braket{ij}_{\alpha}}$  for the corresponding bond. There are as many vectors $\ket{u} \in \mathcal{H}_u$ as bonds in the lattice. The operators $c_i^\alpha$ and $c_j^\alpha$ flip the gauge field ${u}_{\braket{ij}_{\alpha}} \to -{u}_{\braket{ij}_{\alpha}}  $ in a given bond%, see e.g. \eqref{C:c-action}.
\begin{equation}
     c^\gamma_i \ket{u} =  \ket{u'} \; , \quad
    c^\gamma_j \ket{u} = - \im u_{\braket{ij}_\gamma}  \ket{u'} \; ,
\end{equation}
where $ij$ are in a $\gamma$ bond with $i$ in the $\rmA$ sublattice and $u'$ is obtained from $u$ by flipping the sign of the gauge field in the $\braket{ij}_\gamma$ bond.
The state for the translational invariant gauge $u_0=+1$ is defined as the state $\ket{u_0}$ annihilated by  the complex fermion $\frac{1}{2}\left( c_i^\alpha - \im c_j^\alpha \right)$ for each bond $\braket{ij}_\alpha$.

The extended Hilbert space factorizes as the product for gauge field and matter fermions $\mathcal{H}=\mathcal{H}_u \otimes \mathcal{H}_{m}$.  Hhere $\mathcal{H}_u$ is spanned by the basis of eigenvectors of $\hat{u}_{\braket{ij}_\alpha}$ for each of the $3N$ bonds and  $\mathcal{H}_m$ can  be %\footnote{For the matter fermions, there is no-canonical basis, I choose the }
spanned by the basis of eigenvectors of $\im c_i^0 c_{i_z}^0 $ for each one of the $N$ unit cells.
There is a direct connection of the gauge space with the flux of the loop operators in the spin model.
The flux sectors of $W_p$ are linear subspaces of the gauge space $\mathcal{H}_u$ % with $2^{2N}$ elements
and they are related by
\begin{equation}
    W_p = \prod_{\braket{ij}_{\alpha} \in \partial p } \hat{u}_{\braket{ij}_{\alpha} } \; .
\end{equation}
Within each flux sector, we can choose %\footnote{The are $2^{2N-1}$ p}
one out of the $2^{2N}$ gauge configurations\footnote{There are $2N$ gauge transformation $D_i$ in $2N$ sites and $2^{2N}$ gauge transformations of the form $D_g$. %Thus the $2^{3N}$ states in $\mathcal{H}_u$ corresponds to $2^N$ %, and $-1$ since the transformation $\prod_iD_i$ does nothing  -- actually it does in the extended space
.}
, and then the total Hamiltonian $H[u]$ is quadratic in the $c_i^0$. % fermions within the sector of common eigenvectors of $u_{\braket{ij}}$.
Therefore, the Kitaev model is equivalent to a free hopping problem on Majorana fermions $c_i^0$ coupled to a static $\mathbb{Z}_2$-gauge field $\hat{u}_{\braket{ij}_{\alpha} }$.



The choice of gauge field in a given flux configuration should not matter. Let us see how the $\mathbb{Z}_2$ gauge transformation acts on the gauge fields.
The gauge transformations are the product of the operators $\rmD_i$ over a subset $g$ of the lattice,
\begin{equation}
    \rmD_g=\prod_{i\in g} \rmD_i \; .
\end{equation}
The action on any Majorana fermion is $\rmD_g c_i^\mu \rmD_g = \tau_i c_i^\mu$ with $\mathbb{Z}_2$ variable $\tau_i=-$ if $i \in g$  or $\tau_i=+$ if $i \notin g$.
Thus, the action of the gauge transformation in the gauge field operator is
\begin{equation}
     \hat{u}_{\braket{ij}_{\alpha} } \quad \longrightarrow \quad \rmD_{i'} \, \hat{u}_{\braket{ij}_{\alpha} }\, \rmD_{i'} = %\sigma_{ij}\; c_j^\mu \; \quad \sigma_{ij}=
    \left\{\begin{array}{cc}   - \hat{u}_{\braket{ij}_{\alpha} } & i'\in \braket{ij}_{\alpha} \\[6pt]   + \hat{u}_{\braket{ij}_{\alpha} } & i'\notin \braket{ij}_{\alpha} \end{array}\right. \;.
\end{equation}
This implies, reducing to the singleton case $g=\{i\}$ and for the state of gauge field $\ket{u}$, that $\rmD_i\ket{u}\propto \ket{\tilde{u}} $ with gauge field $\tilde{u}$ equal to $u$ but flipped in the three bonds ${\braket{ij}_{\alpha} }$ connected to the site $i$.
More generally, the gauge transformation $\rmD_g$ forms an equivalence class of gauge field $u$, that $u\sim u' \Leftrightarrow \rmD_g \ket{u} \propto \ket{u'}$ for some $g$.
In the flux-free sector, where the flux $W_p=\prod_{\braket{ij}\in \partial p} u_{\braket{ij}}$ is $+1$ on each hexagonal plaquette, there is only one equivalence class for fields in a plane, which are all equivalent to the gauge field $u_0$ which is equal to $(u_0)_{\braket{ij}_\alpha}=+1$ in every oriented bond. In the torus, however, there are four equivalence classes  for gauge fields in the flux-free sector. A choice for them is $u_0$, $u_x$, $u_y$ and $u_z$, where $u_\gamma$ is obtained from $u_0$ by flipping the gauge field along a string of $\gamma$ bonds in the direction $\mb{N}_\gamma$.






The gauge field $\hat{u}_{\braket{ij}_\alpha}$ is defined for each bond. We will use the notation where $\im c_i^\alpha c_j^\alpha$ is ordered with sublattice index $i\in \rmA$ and write
%For $i\in\mathrm A$, we can define the gauge field
\begin{equation}
    \hat{u}_i^\alpha = \im c_i^\alpha c_{i_\alpha}^\alpha  .
\end{equation}
In terms of these gauge operators, the Hamiltonian is
\begin{align}
(K\text{ term}) &= \sum_{i\in \mathrm{A}}\sum_{\alpha}  K_\alpha \im c_i^0 c_i^\alpha \im c_{i_\alpha}^0 c_{i_\alpha}^\alpha
    =-\sum_{i\in\mathrm A}\sum_{\alpha} K_{\alpha} \hat{u}_i^\alpha \, \im c_i^0 c_{i_\alpha}^0 \; .
\end{align}
The same can be done for the remaining of the Hamiltonian using $c^0 c^\gamma = - c^\alpha c^\beta$:
\begin{align}
(\kappa\text{ term})
&=- \kappa
    \sum_{i\in \mathrm{A}} \sum_{\gamma(\beta \alpha )}  \im  c_i^0 c_i^\alpha \im c_{i_{\alpha}}^0 c_{i_{\alpha}}^\gamma  \im c_{i_{\alpha\bar{\beta}}}^0 c_{i_{\alpha\bar{\beta}}}^\beta
    + (\rmA \leftrightarrow \rmB )
    % - \kappa
    %\sum_{i\in \mathrm{B}} \sum_{\gamma(\beta \alpha )}   \im c_i^0  c_i^\alpha  \im c_{i_{\bar{\alpha}}}^0 c_{i_{\bar{\alpha}}}^\gamma  \im c_{i_{\bar{\alpha}{\beta}}}^0 c_{i_{\bar{\alpha}{\beta}}}^\beta
    \;,     \\[5pt]
    &=
    +  \kappa
    \sum_{i\in \mathrm{A}} \sum_{\gamma(\beta \alpha )}  \im  c_i^0  c_{i_{\alpha\bar{\beta}}}^0  \im  c_i^\alpha c_{i_{\alpha}}^\alpha  \im c_{i_{\alpha\bar{\beta}}}^\beta  c_{i_{\alpha}}^\beta
    + (\rmA \leftrightarrow \rmB )
    %+ \kappa
    %\sum_{i\in \mathrm{B}} \sum_{\gamma(\beta \alpha )}   \im c_i c_{i_{\bar{\alpha}{\beta}}} \im c_{i_{\bar{\alpha}}}^\alpha c_i^\alpha    \im  c_{i_{\bar{\alpha}}}^\beta  c_{i_{\bar{\alpha}{\beta}}}^\beta
    \; , \\[5pt]
    &=
    +  \kappa
    \sum_{i\in \mathrm{A}} \sum_{\gamma(\beta \alpha )}  \im  c_i^0  c_{i_{\alpha\bar{\beta}}}^0  \hat{u}_i^\alpha \hat{u}_{i_{\alpha\bar{\beta}}}^\beta
    + (\rmA \leftrightarrow \rmB )
   % + \kappa
   % \sum_{i\in \mathrm{B}} \sum_{\gamma(\beta \alpha )}   \im c_i c_{i_{\bar{\alpha}{\beta}}} u_{i_{\bar{\alpha}}}^\alpha     u_{i_{\bar{\alpha}}}^\beta
    \; .
\end{align}
%Where for $i\in\mathrm B$  then $i_{\bar{\alpha}}\in \mathrm A$ and $u_{i_{\bar{\alpha}}}^\alpha=\im c_{i_{\bar{\alpha}}}^\alpha c_i^\alpha$ (since $i_{\bar{\alpha}\alpha}=i$).
The term $(\rmA \leftrightarrow \rmB )$ follows similarly with a sum under the $\rmB$ sublattice.
To find the final form of the Hamiltonian let us re-index the sum over $i\in\mathrm B$ to a sum over $i\in\rmA$ (that is $i_{\bar{\alpha}} \in \mathrm B$)
\begin{equation}
    H[u]=  \sum_{i\in\mathrm A} \Bigg[ \!
    -\!\sum_{\alpha} K_{\alpha} u_i^\alpha \, \im c_i^0 c_{i_\alpha}^0
    + \kappa \sum_{\gamma(\beta \alpha )} \Big(
    \im  c_i^0  c_{i_{\alpha\bar{\beta}}}^0  u_i^\alpha u_{i_{\alpha\bar{\beta}}}^\beta  +
    \im c_{i_\alpha}^0 c_{i_{ {\beta}}}^0 u_{i}^\alpha     u_{i}^\beta \Big) \Bigg] \; .
    \label{eq:Hu}
\end{equation}
Where, for a fixed gauge configuration $u$, the Hamiltonian \eqref{eq:Hu} is quadratic in the $c^0$ fermions.
The Hamiltonian is block diagonal, $H=\sum_u H[u] \ket{u}\bra{u}$ with blocks $H[u]$.


%In a gauge configuration given by gauge field $\hat{u}_{\braket{ij}_\alpha}$ the Kitaev Hamiltonian can be written


\section{Vaccum state}

%One of the ground states of the Kitaev model is the vacuum state.
The ground states of the Kitaev Hamiltonian are the lowest physical energy states in the flux-free configuration. The simplest configuration is the translational invariant vacuum state.
In this state, the flux through each plaquette is $W_p=+1$ as well as the flux through the non-trivial loops of the torus $W^{(1)}=W^{(2)}=+1$.
%The zero flux sector is interesting because the ground state is known to be on it and there is a translation invariant gauge choice that allows us to diagonalize the single-particle Hamiltonian.
%The ground states are in the sector with no vortices.
A gauge choice for this flux configuration is the homogeneous $u_i^\alpha=+1$. On this gauge, the Hamiltonian becomes
\begin{equation}
    H_0 = H[u_0=+1] =
    \sum_{i\in\mathrm A} \Bigg[ \!
    -\!\sum_{\alpha} K_{\alpha}  \im c_i^0 c_{i_\alpha}^0
    + \kappa \sum_{\gamma(\beta \alpha )} \Big(
    \im  c_i^0  c_{i_{\alpha\bar{\beta}}}^0  +
    \im c_{i_\alpha}^0 c_{i_{ {\beta}}}^0 \big)   \Bigg] \; .
\end{equation}
To diagonalize this translational invariant Hamiltonian, we use the fermions in the momentum space.
Let us define the Fourier transformed modes as
\begin{equation}
    \begin{split}
         c_{\mb{k},\rmA }^0 &= \frac{1}{\sqrt{2N}}\sum_{i\in \rmA } e^{-\im \mb{k}\cdot \mb{R}_i} c_{i_{\phantom{z}}}^0 \; , \quad
    c_{i \in \rmA }^0 = \sqrt{\frac{2}{N}} \sum_{\mb{k}\in{\text{BZ}}} e^{\im \mb{k}\cdot \mb{R}_i} c_{\mb{k},\rmA}^0 \; ,
    \\
    c_{\mb{k},\rmB }^0 &= \frac{1}{\sqrt{2N}}\sum_{i\in \rmA } e^{-\im \mb{k}\cdot \mb{R}_i} c_{i_z}^0 \; , \quad
    c_{i \in \rmB }^0 = \sqrt{\frac{2}{N}} \sum_{\mb{k}\in{\text{BZ}}} e^{\im \mb{k}\cdot (\mb{R}_{i}-\bm{\delta}_z ) }   c_{\mb{k},\rmB}^0 \; ,
     \label{eq:FT}
    \end{split}
\end{equation}
where $\nu=\mathrm{A}, \mathrm{B}$ is the corresponding sublattice of $i$.
For the $B$ sublattice we include the phase factor $e^{-\im \mb{k}\cdot\bm{\delta}_z}$ so that we get the same definition as is \cite{Kitaev2006}\footnote{Here $\mb{R}_i$ is the position vector and not unit cell position.}.
%Note that this transformation differs from the transformation used by Kitaev \cite{Kitaev2006} or in Ref. \cite{Panigrahi_2023}.
%In this transformation, there is a phase difference $e^{\im \mb{k}\cdot\bm{\delta}_z}$ between the fermions in the same unit cell. This will lead to the same \textit{phase} change in the momentum space Bloch Hamiltonian (Kitaev functions $f$ and $\Delta$), but no changes in physical quantities.

The Brillouin zone (BZ) is the surface given by points
\begin{equation}
    \text{BZ} = \left\{ \mb{k}=\frac{k_1}{2\pi}\mb{b}_1 + \frac{k_2}{2\pi}\mb{b}_2 \; , \; k_i=\frac{2\pi}{L_i}n_i, n_i = 0,1,\cdots,L_i-1
    \right\}
\end{equation}
for $\mb{b}_{1/2}= 2\pi \big(\pm1,\frac{1}{\sqrt{3}}\big)$ so that $\mb{a}_{{a}} \cdot \mb{b}_{{b}}=2\pi \delta_{{a}{b}}$, and $L_{1/2}$ are the number of unit cells of the honeycomb lattice along the direction $\mb{a}_{1/2}$. The Brillouin zone is shown in Fig.~\ref{fig:BZ}.
We identify points in the BZ that differ by integer multiples of the basis vectors $\mb{b}_{1/2}$ so that the BZ is a  discretization of the torus with $N=L_1L_2$ points.
I will use $k_1, k_2 \in \left[ 0,2 \pi\right)$  to denote the coordinates of $\mb{k}$ with respect to the basis vectors $\mb{b}_i$ while the components of the vectors are $\mb{k}\doteq (k_x,k_y)=(k_1-k_2,\frac{k_1+k_2}{\sqrt{3}})$. For $\mb{k}, \mb{k}' \in {\text{BZ}}$ holds the identity $\sum_{i\in\mathrm A} e^{\im (\mb{k}-\mb{k}' )\cdot \mb{R}_i}=N \delta_{\mb{k},\mb{k}'}$.

The Hermiticity of the Majorana fermions $c_i$ implies
\begin{equation}
    c_{\mb{k},\nu}^\dagger=c_{-\mb{k},\nu }^\pd \; ,
\end{equation}
which relates complex fermion operators for antipodal points $\mb{k} \sim  -\mb{k}$. This relation ensures that we actually have $N$ complex fermions\footnote{Which is the correct number of fermionic degrees of freedom for a system of $2N$ Majorana fermions $c_{i}$.} $c_{\mb{k},\nu}$ as long as we restrict $\mb{k}$ to take values over the $N/2$ points of the half-Brillouin zone.
The half-BZ is obtained from the BZ by removing points so that no antipodal points are in it\footnote{To avoid problems, take $N$ to be even and shift the BZ to not include $\mb{0}$, so that we have exactly $N$ complex fermions.}.
From the anticommutation relation for $c_i$, follow that the complex fermions satisfy the usual anticommutation relation for $N$ independent fermions
\begin{equation}
    \{ c_{\mb{k},\nu }^\dagger,c_{\mb{k}',\nu' }^\pd  \}=\delta_{\mb{k},\mb{k}'} \; , \quad
    \{ c_{\mb{k},\nu },c_{\mb{k}',\nu' }^\pd  \}= 0 \; ,
\end{equation}
for $\mb{k},\mb{k}'$ in the half-BZ.
Note that we can use the fermions $c_{\mb{k},\nu}$ in the full BZ, but we need to be careful to use
\begin{equation}
     \{ c_{\mb{k},\nu },c_{\mb{k}',\nu' }^\pd  \}= \delta_{\mb{k},-\mb{k}'} \; , \quad \mb{k},\mb{k}' \in \text{BZ} \; .
\end{equation}
As we will see, when expressing in terms of the normal modes that diagonalize the Hamiltonian, we can trade off the restriction of $c_{\mb{k},\nu}$ to the half-BZ for a restriction of the normal modes to non-negative energy $\varepsilon(\mb{k})\geq 0$ in the full BZ.

% which is the surface obtained by the Torus of the Brillouin zone identifying antipodal points.
%The particular choice for the values of $\mb{k}$ in the half-Brillouin-zone is not important, but to get exactly $N/2$.



%\begin{marginfigure}
\begin{figure}
    \centering
    \includegraphics[width=.4\linewidth]{Figures/2/BZ.pdf}\hspace{5mm}%
    \includegraphics[width=.4\linewidth]{Figures/2/half-BZ.pdf}
    \caption{
    Two geometries for the Brillouin zone (BZ, in pink)  and the half-BZ (in gray). Points that differ by integer multiples of $\mb{b}_{1/2}$ are equivalent. In particular, opposite edges of the BZ are identified. Points of high symmetry $\bm{\Gamma},\mathrm{K},\mathrm{K}',\mathrm{M}$ are highlighted.   }
    \label{fig:BZ}
\end{figure}



The Hamiltonian in terms of the Fourier transformed modes is
\begin{align}
    H_0 &= \frac{1}{2} \sum_{\mb{k}} \Big[
    2\gamma(\mb{k}) c_{\mb{k}\mathrm{A}}^\dagger c_{\mb{k}\mathrm{B}}^\pd +
     c_{\mb{k}\mathrm{A}}^\dagger c_{\mb{k}\mathrm{A}}^\pd \, \Delta(\mb{k})+
      c_{\mb{k}\mathrm{B}}^\dagger c_{\mb{k}\mathrm{B}}^\pd \, \Delta(-\mb{k}) \Big]
\end{align}
with functions
\begin{align}
        \gamma(\mb{k}) &=
        -2 \im  \sum_\alpha K_\alpha e^{\im \mb{k}\cdot (\mb{\delta}_\alpha-\mb{\delta}_z)}
        = - 2 \im \Big[  K_x e^{-\im \mb{k}\cdot \mb{a}_1} +K_y e^{-\im \mb{k}\cdot \mb{a}_2} + K_z  \Big] \, ,
        \\
    \Delta(\mb{k}) &=
    \text{Re}\Big[  4 \im \kappa \sum_{\gamma(\alpha\beta) } e^{\im \mb{k}\cdot(\delta_\alpha-\delta_\beta)} \Big] =
    4 \kappa \sum_\gamma \sin \big( \mb{k}\cdot \mb{N}_\gamma  \big)
\end{align}
where $\mathbf{N}_{\gamma}=\mathbf{N}_{\gamma(\beta \alpha )}=\bm{\delta}_{\beta}-\bm{\delta}_{\alpha}$ so that $\mathbf{N}_{x}=\mb{a}_2, \mathbf{N}_{y}=-\mb{a}_1, \mathbf{N}_{z}=\mb{a}_1-\mb{a}_2$. %The real part in $\Delta(\mb{k})$ appears as we can use the (anti-)symmetry  $\mb{k}\to-\mb{k}$ of the sum to cancel the imaginary part (which is even in $\mb{k}$).
These functions satisfy the relation under momentum reversal
\begin{equation}
     \gamma(\mb{k})^\ast=-\gamma(-\mb{k}) \; , \quad  \Delta(-\mb{k})=-\Delta(\mb{k}) \; .
\end{equation}
Using this symmetry, we can rewrite the Hamiltonian as  % on the the terms as $\sum_{\mb{k}} \gamma(\mb{k}) c_{\mb{k}\mathrm{A}}^\dagger c_{\mb{k}\mathrm{B}}^\pd = \sum_{\mb{k}} \gamma(\mb{k})^\ast c_{\mb{k}\mathrm{B}}^\dagger c_{\mb{k}\mathrm{A}}^\pd$ so that
\begin{align}
    H_0 &=   \sum_{\mb{k}}^{\prime}
    \left(\begin{array}{c} c_{\mb{k}\mathrm{A}}^\dagger \\ c_{\mb{k}\mathrm{B}}^\dagger
    \end{array}\right)^{T}
    \left(\begin{array}{cc}
        \Delta(\mb{k}) & \gamma(\mb{k}) \\
       \gamma(\mb{k})^\ast  & -\Delta(\mb{k})
    \end{array}\right)
    \left(\begin{array}{c} c_{\mb{k}\mathrm{A}}^\pd \\ c_{\mb{k}\mathrm{B}}^\pd
    \end{array}\right)
    =  \sum_{\mb{k}}^{\prime}  c_{\mb{k}}^\dagger h(\mb{k}) c_{\mb{k}}^\pd \; .
\end{align}
with $\sum^{\prime}$ denoting the sum over the half-Brillouin-zone and the block matrix for the Hamiltonian in momentum space is
\begin{equation}
    h(\mb{k}) = \mb{d}(\mb{k})\cdot \bm{\sigma} \; , \quad
    \mb{d}(\mb{k}) = \big(\text{Re}\gamma(\mb{k}), -\text{Im}\gamma(\mb{k}), \Delta(\mb{k}) \big) \; .
\end{equation}
The $x/y$ components of $\mb{d}$ are such that $d_x - \im d_y=\gamma$. This vector satisfies the transformation relation $\mb{d}(-\mb{k}) = \big(-d_x(\mb{k}),d_y(\mb{k}),-d_z(\mb{k}) \big)$ which shows the C-symmetry of the Hamiltonian
\begin{equation}
    h(-\mb{k})^T = - h(\mb{k}) \; .
\end{equation}



We can diagonalize the Hamiltonian by a unitary transformation $h(\mb{k})=U(\mb{k}) \left(\begin{smallmatrix} \epsilon(\mb{k}) & 0 \\ 0 & -\epsilon(\mb{k})
\end{smallmatrix}\right) U^\dagger (\mb{k})$ in which the single-particle energies are $\pm\epsilon(\mb{k})=\pm|\mb{d}(\mb{k})|$. The corresponding eigenmodes $\alpha_{\mb{k} }$ are given by
\begin{equation}
    \left(\begin{array}{c} \alpha_{\phantom{-}\mb{k}}^\pd \\ \alpha^\dagger_{-\mb{k}}
    \end{array}\right) = U^\dagger (\mb{k}) \;
    \left(\begin{array}{c} c_{\mb{k}\mathrm{A}}^\pd \\ c_{\mb{k}\mathrm{B}}^\pd
    \end{array}\right) \; , \quad
    U =
    \frac{1}{\sqrt{2\epsilon(\epsilon\!+\!\Delta)}}\, \left( \begin{array}{cc}
         \!\!\epsilon + \Delta\!\! & -\gamma \\
         \gamma^\ast & \epsilon +\Delta\!\!
    \end{array} \right)
\end{equation}
where $\epsilon(\mb{k})=|\mb{d}(\mb{k})|=-\epsilon(-\mb{k})$ %and $\epsilon(-\mb{k})=-|\mb{d}(\mb{k})|$
and $\mb{k}$ is in the half-Brillouin-zone. The reason to write these complex fermions with the same label, $\alpha_{\phantom{-}\!\!\mb{k}}^\pd$ and $\alpha^\dagger_{-\mb{k}}$, comes from the charge conjugation symmetry %\footnote{Some people call it particle-hole symmetry, and here I rephase ``there is no need to become philosophical about whether to call it charge conjugation or particle–holes symmetry''.}
of the Majorana fermion Hamiltonian %($h(\mb{k})=-\sigma_x h^{T}(-\mb{k}) \sigma_x$)
together with $c_{\mb{k}\nu}^\dagger=c_{-\mb{k}\nu}^\pd\!$ and $U(-\mb{k})=-U^\ast (\mb{k})$.
In terms of the normal modes, the Hamiltonian in canonical form is
\begin{equation}
    H_0 =  \sum_{\mb{k}}^{\prime} \epsilon(\mb{k}) \, \big( \alpha^\dagger_{\mb{k}}\alpha_{\mb{k}}^\pd - \alpha_{-\mb{k}}^\pd\alpha^\dagger_{-\mb{k}} \big)
    =  \sum_{\mb{k}} |\mb{d}(\mb{k}) | \, \alpha^\dagger_{\mb{k}}\alpha_{\mb{k}}^\pd + E_0 \label{eq:k-h0}
\end{equation}
and the ground state energy is $E_0=-\sum_{\mb{k}}^{\prime} \,  |\mb{d}(\mb{k})|=-\frac{1}{2} \sum_{\mb{k}} \,  |\mb{d}(\mb{k})|$.
Thus, the ground state of $H_0$ is the state $\ket{\varnothing, u_0 }$,
\begin{equation}
    \forall \mb{k} \quad \alpha_{\mb k}^\pd \ket{\varnothing, u_0 } = 0 \; ,
\end{equation}
annihilated by all $\alpha_{\mb{k}}$ with $\mb{k}$ in the full Brillouin zone. % (here we only have quasi-particles with positive energy).
The excited states are
\begin{equation}
    \ket{\Psi, u_0 } = \prod_{\mb k}  \left(\alpha_{\mb k}^\dagger \right)^{n_\mb{k}} \ket{\varnothing, u_0 }
\end{equation}
for each $n_{\mb k} \in \{0,1\}$ and the energy of these states in $H_0$ follows from \eqref{eq:k-h0}.


The quasi-particle energy dispersion is
\begin{equation}
    \epsilon(\mb{k})
    = |\mb{d}(\mb{k})|
    = \sqrt{ \, |\gamma(\mb k)|^2 + \Delta^2(\mb k) \, } \; .
\end{equation}
The low energy states for $\kappa=0$ are in a gapless phase when there are real solutions to $K_x e^{\im k_1}+K_y e^{\im k_2}+K_z = 0$, which happens when the couplings satisfy the triangle inequalities $|K_\alpha|+|K_\beta|\geq |K_\gamma|$ for $(\alpha\beta\gamma)$ different and is gapped otherwise. % In particular, we found a gapless with conic dispersion.% in the isotropic limit $K_\alpha=K$.
We are mostly interested in the phase near the isotropic limit $K_\alpha=K$.
This phase becomes gapped for non-zero $\kappa$.

%\reminder{ Zero spin expectation value, spin-spin correlation ultra-local, Majorana-majorana long-ranged for $\kappa=0$  }



Consider the low-energy effective theory in the isotropic  $K_\alpha=K$ limit and small non-zero $\kappa$.
The single particle energy dispersion has minima in the $\mathrm{K}$ and $\mathrm{K}'$ points, with are at $\pm\mb{k}_{\ast}$ with $\mb{k}_{\mathrm{K}}=\frac{2}{3}\mb{b}_1+\frac{1}{3}\mb{b}_2$. The conic dispersion is now hyperbolic.
Considering up to linear order deviations in momentum close to the $\mathrm K$ point,  $\mb{k}=\mb{k}_{\mathrm{K}} + k_1 \mb{b}_1+ k_2 \mb{b}_2 $,
we have
\begin{equation}
    \begin{split}
        \gamma(\mb{k}) &\approx
        -2 K\big[ k_1 e^{-2\pi \im \frac{2}{3}}+k_2 e^{-2\pi \im \frac{1}{3}}\big]
        %= - 2 K \Big[ (k_1+k_2) \frac{1}{2}+\im (k_1-k_2)\frac{\sqrt{3}}{2} \Big]
        =-\im v \big(   k_x -  \im k_y \big) \; , \\[5pt]
        \Delta(\mb{k})
        &\approx \Delta_0
        =  4 \kappa \Big[ -\sin(2\pi \frac{2}{3})+ \sin(2\pi \frac{1}{3})+\sin(2\pi \frac{1}{3}) \Big]
        =6\sqrt{3} \kappa
    \end{split}
\end{equation}
for $v  =\sqrt{3} K, (k_x,k_y)=(k_1-k_2,\frac{k_1+k_2}{\sqrt{3}})$ and $\Delta_0=\Delta(\mb{k}_{\mathrm{K}}) $
which simplify $\mb{d}(\mb{k}) = (\mathrm{Re}\gamma,-\mathrm{Im}\gamma,\Delta)\approx ( v k_y,  v k_x,\Delta_0)$ giving the energy dispersion $\approx \sqrt{ \, v^2 (k_x^2+k_y^2) + \Delta^2_0 \; }$.


For $\kappa \neq 0$, this phase becomes gapped. Kitaev demonstrated \cite{Kitaev2006} that this gapped phase possesses non-trivial topological properties characterized by a non-zero Chern number $=\text{sgn}(\kappa)$ in the phase that is gapless at $\kappa=0$. Analogous to the integer quantum Hall effect, the non-zero Chern number implies the presence of anyonic properties in vortices.
Due to the property of the anyons, this phase is known as the non-Abelian phase of the Kitaev model, in contrast with the Abelian gapped phase for $\kappa=0$ and $|K_\gamma|>|K_\alpha|+|K_\beta|$.
Moreover, the system must have chiral edge modes in an open geometry.
These chiral edge modes are Majorana fermions, thus the edge modes do not contribute to the electrical current as in the conventional Quantum Hall Effect. Instead, they are expected to contribute to a quantized thermal energy current along the edge.



\section{Physical states}

The connection between a state $\ket{\Psi ; u }\in \mathcal{H}_f$ for the single-particle Hamiltonian $H[u]$ and the state for the many-body Hamiltonian requires a gauge projection to ensure $\rmD_i=+1$. By taking the product with the corresponding gauge state in $\mathcal{H}_u$ and performing the projection, we find the state
\begin{equation}
    \ket{\Psi} =  \prod_{i }\left( \frac{1+\rmD_i}{2} \right) \ket{u}\otimes\ket{\Psi ; u }
\end{equation}
which, if the resulting state is non-zero,  is a state in the physical subspace $\rmD_i\ket{\Psi}=+\ket{\Psi}$ and an eigenstate of $H$ with the same energy as the eigenvalue of  $\ket{\Psi ; u }$ in $H[u]$. To find ground states, we need to understand how the projection operator acts in the gauge and matter sector to, for instance, understand when we have a non-zero resulting state.



We define the projector operators as
\begin{equation}
    \cP = \prod_{i\in \mathcal{L} } \frac{1+\rmD_i}{2} = \frac{1}{2^{2N}}  \sum_{g \in \mathrm{P}(\mathcal{L})} \rmD_g \; , \quad D_g = \prod_{i\in g } D_i \; ,
\end{equation}
which project out the states with $\rmD_i=-1$.
Here $\mathcal{L}$ is the set of all $2N$ honeycomb lattice points and $\mathrm{P}(\mathcal{L})$ is the power-set containing all subsets $g$ of $\mathcal{L}$.
There are $2^{2N}$ terms in the sum on the right, including the terms corresponding to the empty subset $1= \prod_{i\in \varnothing } D_i$ and to the whole set
\begin{equation}
    \rmD:= \prod_{i\in \mathcal{L} } \rmD_i = (-1)^{\hat{n}_c}\; ,
\end{equation}
where $(-1)^{\hat{n}_c} = \prod_{i,\mu} c_i^\mu$ is the total fermionic parity. %where the parity for two Majorana fermions is equal to  $(-1)^{\hat{n}_c}= \im c^\prime c^{\prime\prime}$. %, and similarly for any even quantity\footnote{It is independent of how we pair the Majorana fermions.}.
Note that the gauge transformation $\rmD_{g}$ and $\rmD_{\mathcal{L}\backslash g}$ are not independent since $\rmD_{\mathcal{L}\backslash g} = \rmD \; \rmD_g$ and $\rmD$ performs a gauge transformation on the whole lattice\footnote{which is $+1$ if the state is not annihilated by $\cP$.}.
A better way to express the projection operator is to factorize it as \cite{Pedrocchi2011,Yao_2009,Yao_2010}
\begin{equation}
    \cP =  \left(
    \frac{1}{2^{2N-1}}  \sum_{g \in \mathrm{P}(\mathcal{L})}^\prime  \rmD_g
    \right) \; \left( \frac{1+\rmD}{2} \right)
    = \mathcal{S} \;   \left( \frac{1+\rmD}{2} \right)
    \; ,
\end{equation}
where $\mathcal{S}$ consists of a symmetrization over all independent gauge transformations and $ \left( \frac{1+\rmD}{2} \right) $ projects in the state with even total fermionic parity $(-1)^{\hat{n}_c}=+1$.


\begin{comment}
The ground state is obtained by applying $\mathcal{S}$ to the lowest energy state with even fermionic parity.
It is convenient to separate the parity of the $c_i^0$ fermions from the $c^{1,2,3}_i$ ones, which can be collected using the gauge field operators $u_i^\alpha = \im c_i^\alpha c_{i_\alpha}^\alpha$ to find \cite{Pedrocchi2011}
\begin{equation}
    \rmD = \pi_{c^0} \; %(-1)^{L_1+L_2}
    \prod_{i,\alpha} u_i^\alpha \;
\end{equation}
for periodic boundary conditions with $L_1+L_2$ even and $\pi_{c^0}$ is the parity of the $c^0$ fermions.

Define gauge transformations $\rmD_{L_1}$ and $\rmD_{L_2}$,
\begin{equation}
    \rmD_{L_{1/2}} = \prod_{i\in\ell_{1/2}} \rmD_i \; ,
\end{equation}
along closed loops $\ell_{1/2}$  of the torus. Since the loops are closed, this transformation does

\end{comment}

%\reminder{.. degeneracy of the ground state, entanglement entropy}

%\section{Excited states}

%\reminder{.. general expression (like Bogoliubov, Thouless theorem) for a given configuration, gauge configurations,  total number of states}

Using this projector we can find all eigenstates of the Kitaev Hamiltonian.
In each flux sector, we choose a gauge field.
This gives us a single-particle Hamiltonian $H[u]$. By diagonalizing this $2N$ square Hamiltonian we find a vacuum $\ket{\varnothing;u}$ and quasi-particle index $k=1, \cdots , N$  independent normal-modes $\alpha_k(u)$.
These modes are such that
\begin{equation}
    H[u]
    =  \sum_{k} \varepsilon_k(u) \, \alpha^\dagger_{{k}}(u)\alpha_{k}^\pd(u) + E_0[u] \; , \quad
    E_0(u)=- \frac{1}{2} \sum_{k} \varepsilon_k(u)
\end{equation}
With this procedure, any state is indexed by the flux sector, and the collection of occupation numbers ($n_k=0$ or $1$) of the normal modes
\begin{equation}
    \ket{\Psi_{ \{n_k\},W }} =  \left( \frac{1+\rmD}{2} \right) \;  \left[ \mathcal{S} \ket{u}\otimes \; \prod_{k=1}^{N} \left(\alpha^{\dagger}_{k}(u) \right)^{n_k}\ket{\varnothing;u} \right] \; . \label{eq:2-physical-state}
\end{equation}
The symmetrization $\mathcal{S}$ ensures that the state does not depend upon the specific gauge choice $u$. Naively, for $N$ plaquettes, there are $2^{N}$ possible flux sectors and for $N$ normal modes there are $2^{N}$ possible states, which reflect the total of $2^{2N}$ physical spin states.
Actually,  we can only create vortices in pairs due to the constraint $\prod_p W_p=+1$, so we only have  $2^{N-1}$ possible flux sectors, but one needs to consider the flux through the non-trivial Wilson loops $W^{(1)}$ and $W^{(2)}$, giving $2^{N+1}$ possible flux configurations. The parity constraint imposed by $\left( \frac{1+\rmD}{2} \right)$ allows only half of the fermionic states,  thus we can recover all the $2^{N+1}2^{N-1}=2^{2N}$ physical states from \eqref{eq:2-physical-state}.

The energy of the state is independent of the symmetrization operator $\mathcal{S}$.
Meaning that the eigenvalue of $\ket{\Psi,u}$ in $H[u]$ is the same as the eigenvalue of $\mathcal{S}\ket{u}\otimes \ket{\Psi,u}$ in $H$, if the total parity is even.
%Furthermore, to compute the expectation of operator $\mathcal{O}$ that preserves the
Thus, the symmetrization operator is not relevant in the intermediate calculations, % but to ensure the total even fermionic parity it is.
but it is crucial to ensure that the total fermionic parity is even.
For a given flux sector, the lowest energy state in it is given by either the ground state or the first excited state of $H[u]$ for any gauge $u$ corresponding to this flux sector.
Which state is the correct physical lowest energy state depends on the fermionic parity of the ground state $\ket{\varnothing,u}$ of $H[u]$, and it is
\begin{equation}
    E_{0}[W]= \left\{ \begin{array}{cl}
        E_{0}[u] & \; , \quad \text{if} \quad \rmD\ket{\varnothing,u}= + \ket{\varnothing,u}  \\[8pt]
        E_{0}[u] + \varepsilon_1(u) &\; , \quad \text{else} \; .
    \end{array} \right. \label{eq:gs-en}
\end{equation}
The state corresponding to the latter is $\alpha_1^\dagger(u) \ket{\varnothing,u}  $.
To find the many-body ground states of $H$ one needs to compute \eqref{eq:gs-en} for all different flux configurations and find the one that minimizes it. As discussed previously, by Lieb's theorem applied to the honeycomb lattice with periodic boundary conditions, one only needs to consider four different flux configurations corresponding to the flux-free $W_p=+1$ and $W^{(1/2)}=\pm 1$.


%\section{Ground states }


\section[{$\mathbb{Z}_2$}-flux and Majorana modes]{Visons and Majorana modes}

In the Kitaev model, the spin degrees of freedom are mapped into free Majorana fermions $c^0_i$ coupled to a static background gauge field $u_{\braket{ij}_\gamma}$. We call it a gauge field as it enters in the hopping Hamiltonian similar to the Peierls phase factor $e^{\im A_{ij} }$ for lattice complex coupled to $U(1)$ gauge.
The excitations in this background field are the $\mathbb{Z}_2$ vortices. From the point of view of the Majorana fermions in the background of a single vortex, it changes the sign of the hopping in a string from the center of the vortex to a point in the infinity, see Fig. \ref{fig:2-vortex-string}.
\begin{figure}
    \centering
    \includegraphics[width=.5\linewidth]{Figures/2/1vortex.pdf}
    \caption{A configuration with one isolated vortex, represented in pink with flux $\pi$, line extent to infinity in an infinite open planar geometry. The gauge field $u_{\left<ij\right>_\gamma}$ is $+1$ for every bond except along the string, which has value $-1$. The Majorana fermions get a $\pi$ phase by encircling the vortex, similar to how electrons acquire an Aharonov-Bohm phase by encircling a magnetic flux.
    }
    \label{fig:2-vortex-string}
\end{figure}
This change in sign, in the lattice fermionic language, is equivalent to introducing anti-periodic boundary condition to the Majorana fermions, making then acquire a $\pi$ phase ($-1$ sign) as they encircle the vortex. This is in complete analogy to the Aharonov-Bohm phase $\frac{e}{\hbar}\Phi$ obtained by the electron by encircling a magnetic flux $\Phi$, with the difference that this phase (flux) is not continuous but rather discrete, $0$ or $\pi$. In this sense, the $\mathbb{Z}_2-$vortex is a $\pi$ flux. The flux is not made of a magnetic field, but of the abstract $Z_2$ field to which the Majorana fermions are coupled.


We are interested in a particular class of excited states with a low density of vortices.
A configuration with vortices can be obtained for example from the flux-free standard gauge by flipping the gauge field along one bond $\braket{ij}_\gamma$. This produces a pair of vortices $W_p=-1$ for the two plaquettes $p$ that share this bond. We can separate these vortices with a gauge string, which is to flip the gauge field along some continuous sequence of bonds. The effect is to move the vortices, which will be found on both ends of the gauge string.

We cannot use translational invariance to find the ground state analytically for configuration with vortices. Still, we find the ground state by numerically diagonalizing the Hamiltonian $H[u]$ in \eqref{eq:Hu}. The ground state for the gauge configuration with two isolated vortices is about $0.04 \, |K|$ more energetic than for the configuration with adjacent vortices in the thermodynamic limit see table \ref{tab:I-1-en}. Note that the energies $E_0$ and $E_{2\text{v}}$ for the configuration with and without vortices are extensive with the number of site $N$, but its difference $\Delta_{2\text{v}}=E_{2\text{v}}-E_0$ is in independent of $N$ for large $N$.

Apart from the $\kappa$ term\footnote{We are still using $\kappa\neq0$ in the bulk, I am only ignoring the $\kappa$ term in the flipped bonds for simplicity. It is important to have $\kappa\neq0$ in $H[u_0]$.}, the Hamiltonian is linear in the gauge field. The Hamiltonian for the gauge field $\tilde{u}$ with vortices is related to the Hamiltonian $H[u_0]$ without vortices by
\begin{equation}
    H[\tilde{u}] = H[u_0] + H[\tilde{u}-u_0]\Big|_{\kappa=0} + ...
\end{equation}
where $\tilde{u}-u_0$ is $-2$ along the bonds crossed by the gauge string and it is $0$ elsewhere.
For simplicity take the string along $z$-bond along the $(1,0)$ direction with one vortex at the origin and the other at infinity.
In the continuum limit, which is a low-energy expansion about the Dirac point of the translational invariant Hamiltonian $H[u_0]$,  the first term is described by a 2d Dirac Hamiltonian with mass term $m\sim \kappa$ while $H[\tilde{u}-u_0]$ gives the potential $2 K_z \theta(r_1)\delta(r_2)= 2K_z \frac{\delta(\theta)}{r}$. This problem is equivalent to the 2d Dirac equation with a branch-cut at the line $\theta=0$ that changes the boundary condition by giving a $\pi$ phase shift to the fermions across the string. The solutions of this equation show that there is a single zero energy mode centered at the origin \cite{Sachdev_2023}.
The solution is a Majorana mode and its wavefunction is radially symmetrical and centered on the vortex.
The same conclusion was found without the continuum limit approximation by Kitaev in \cite{Kitaev2006}, see appendix $C$, by showing that the fermionic parity of the problem of a single vortex is odd. This implies an unpaired Majorana mode. From the theory of Majorana fermions, a single Majorana fermion, by particle-hole symmetry, must have zero energy.


%Consider the case of two vortices created by a configuration where $u=1$ everywhere but along a line with $d$ flipped bonds where $u_i^z=-1$. By numerically diagonalizing the Hamiltonian with finite $\kappa=0$ we find the modes $\alpha_k$ to be gapped but for two zero-energy modes\footnote{The energy decays to zero exponentially with the distance of the vortices.}


%In the last section, we saw the excited states on top of a $\mathbb{Z}_2$-flux free sector. Within this sector, we already can conclude the presence of Majorana edge modes in the boundary of the sample in the $B$ phase due to a non-zero topological invariant, the Chern number. We will not get much in detail about these excitations, but they are currently in dispute if they have or have not been observed in thermal Hall conductance experiments \cite{Kasahara2018, Motome2020, Yokoi_2021}.

%The excited states that we will be interested in are in the sector with the $\mathbb{Z}_2$-flux. These states are spanned by the vectors $\mathcal{P}\, \ket{\mathcal{G}} \, \ket{\varphi_{\lambda}(\mathcal{G})\, }$ where $\mathcal{G}$ is a given fixed gauge configuration and $\ket{\varphi_{\lambda}(\mathcal{G})\, }$ are the eigenvectors of matter fermion Hamiltonian\vspace*{-3mm} \begin{equation} H(\mathcal{G}) = \sum_{\,i,j } \frac{\im }{ 4} \, A_{jk}(\mathcal{G}) \; c_j c_k \; . \end{equation} The process to find the desired states follows from fixing a gauge-field $u_{\braket{ij}}$, with more than one $W_p=-1$ corresponding to the fluxes excitations, then calculating the eigenvector from the matrix $A_{jk}(\mathcal{G})$ and select the states with parity $\prod_j D_j = +1$.


\begin{table}[t]
    \centering
    \begin{tabular}{|c||c|c|c|}\hline
       &  Flux-free & Two adjacent vortices \cite{Panigrahi_2023} & Two isolated vortices %\cite{Kitaev2006}
       \\
        \hline
       Energy$/|K|$    %   & $-1.5746$ &  $-1.3074$ & $-1.2674$
       & $E_0=-1.574 \, N$ &  $E-E_0=0.263$ & $E-E_0=0.307$ \\ \hline
    \end{tabular}
    \caption{$E$ is the many-body ground state energy of $H[u]$ with  $\kappa=0$ for three different flux sectors. }
    \label{tab:I-1-en}
\end{table}





In summary,  in the non-Abelian phase, there is a Majorana zero mode (MZM) in each isolated vortex.
These excitations are bound to the vortices.
The bound state from $\mathbb{Z}_2$ fluxes and MZM forms a composite particle %which has a non-Abelian statistic of Ising anyons
$\sigma$ \cite{Kitaev2006}.
Due to the presence of the MZM, the $\sigma$ particles satisfy non-Abelian anyonic mutual statistic.


%\reminder{say we have a complex fermionic mode splitted non-locally into two. The energy difference of these two modes is zero for isolated vortices. When approximating them the levels separate. It can be either occupied or not. From my definition the quasi energies are always positive, but if we let the energy to smoothly change across the parity changing points the   }

\begin{comment}

When we bring the vortices to a finite distance $d$, the Majorana modes acquire finite energy.
The splitting energy  $\varepsilon(d)$ decays exponentially and oscillates with $d$, in Ref. \cite{Lahtinen_2012} they found two zero modes split into two levels of finite energy $E_\psi$ and $E_1$, its difference is $\varepsilon(d)\sim e^{-d/\xi}$ where\footnote{It is defined as $\xi=|\mathrm{Im}(\mb{k})|^{-1}$ for the complex solution of $\varepsilon(\mb{k})=0$.} $\xi=3|K|/\Delta_0$ is the correlation length of the non-Abelian phase. These levels are particle-hole symmetric. % The correct sign for $\varepsilon(d)$ is determined by the parity constraint from choosing the physical state.
In Ref. \cite{Otten_2019}, they found the energy splitting for large distance $\mb{d}$ of two vortices in the continuum limit to be
\begin{equation}
    \varepsilon(\mb{d}) = 8 v \, |\cos (\mb{k}_{\mathrm{K}}\cdot \mb{d} )| \, \sqrt{\frac{1}{\; 2 \pi  d \, \xi} \,} \; e^{- d/\xi}
\end{equation}
Which agrees with \cite{Lahtinen_2011,Lahtinen_2012}.
\end{comment}

\begin{comment}

Consider a configuration with an even number of isolated vortices.
The pair of Majorana modes in each vortex form a complex fermion, say $\psi$, which can be occupied or not.
%If the fermionic state is unoccupied it can be represented by $\ket{{\!\,}_1^{\sigma\sigma}}$ and it is occupied by $\ket{{\!\,}_\psi^{\sigma\sigma}}$.
For even fermionic parity of this two-MZM subsystem, the corresponding Ising anyons are in the sector where their fusions combine into a trivial state with no excitations, $\sigma \times \sigma \to 1$, while for odd parity these anyons have an occupied fermionic state and they fuse into a state with a fermion,
$\sigma \times \sigma \to \psi$.
%Before fusing, the subsystem of two isolated anyons are in either %$\ket{{\!\,}_1^{\sigma\sigma}}$ $\ket{{\sigma\sigma}\atop{1}}$ or $\ket{{\sigma\sigma}\atop{\psi}}$ fusion state, where %upper index indicates the particle type before the collision and lower indices after the collision
The fusion rules are
\begin{equation}
    \sigma \times \sigma = 1 \otimes \psi \; , \quad
    \psi \times \psi = 1 \; , \quad
    \psi \times \sigma = \sigma \; ,
\end{equation}
where $1$, $\sigma$ and $\psi$ represent the superselection sectors of the vacuum, the Ising anyon (vortex carrying an unpaired MZM), and a complex fermion.


\end{comment}

%\reminder{conclude saying the the interaction between vortices is weak -- make bridge to say that the quadrupole is stronger}



\begin{comment}

\section{Non-Abelian Ising anyons}


\end{comment}




\begin{comment}
\reminder{
\begin{itemize}
    \item existence of MZM for each vortex, how the energy is not zero for vortices at finite distance (exponential),
    \item continuum limit
    \item relation to fusion rules,
    \item non-Abelianess,
\end{itemize}
}

\section{Observables}

\reminder{
\begin{itemize}
    \item static correlations in the ground state .. how it is ultra-local.
    \item how does it change in the presence of vortices (it is still local)
    \item can one find an expression in the continuum limit?
    \item relation to experimental tests?
\end{itemize}
}



\end{comment}





